\section{Introduction}

Breadth-first search (BFS) is the typical example of the graph traversal dwarf. It is the
main kernel of Graph500~\cite{graph500} and one of the six kernels of the GAP Benchmark Suite
(\gapbs)~\cite{beamer2015gap}. BFS does almost no arithmetic. Its speed depends on how fast
it can get data from memory, and most of these accesses are random. Because of this, a
model of BFS has to describe data movement, latency and synchronisation, and not
floating-point operations.

The goal of the project was not to write the fastest BFS, but to go through the full
performance engineering process with the \gapbs{} reference code:
\begin{enumerate}
  \item analyse the program with more than one tool, and check that the tools give correct
        results,
  \item explain the measurements with a model whose parameters we measured ourselves,
  \item choose, implement and test optimisations based on the analysis and the model.
\end{enumerate}

In our work plan we wanted to use a laptop and the HAICGU cluster. We did not get access to
HAICGU, so we used the Hábrók cluster of the university instead. We also did more
memory and OpenMP optimisations than planned, because the analysis showed that these parts
take most of the time, and fewer compiler optimisations.
